\documentclass[11pt, a4paper]{article}
% Version: v3-EN
\usepackage[utf8]{inputenc}
\usepackage[T1]{fontenc}
\usepackage[margin=1in]{geometry}
\usepackage{setspace}
\usepackage{lineno}
\usepackage{amsmath}
\usepackage{graphicx}
\usepackage{booktabs}
\usepackage{tabularx}
\usepackage{float}
\usepackage{placeins}
\IfFileExists{garamondx.sty}{\usepackage{garamondx}}{\usepackage{mathptmx}}
\usepackage{caption}
\usepackage{parskip}
\usepackage[colorlinks=true,
            linkcolor=black,
            citecolor=black,
            urlcolor=blue]{hyperref}

% BibTeX için natbib paketi
\usepackage[numbers, super, sort&compress]{natbib}

\captionsetup{labelfont=bf, labelsep=period, font=small}
\doublespacing
\setlength{\emergencystretch}{2em}
\linenumbers

\newcommand{\Fig}[1]{\textbf{Figure~\ref{#1}}}
\newcommand{\Tab}[1]{\textbf{Table~\ref{#1}}}

% ============================================================

\title{\textbf{LINE-1 Depletion in Promoters of
Neurodevelopmental Disorder Genes:\\
Evidence from Genome-Wide Analysis}}

\author{
Can Sevilmi\c{s} \\[4pt]
\small Department of Molecular Biology and Genetics,
Bah\c{c}e\c{s}ehir University, Istanbul, Turkey \\[2pt]
\small \href{mailto:can.sevilmis@bahcesehir.edu.tr}
{can.sevilmis@bahcesehir.edu.tr} \\[2pt]
\small ORCID: \href{https://orcid.org/0000-0002-9180-1924}{0000-0002-9180-1924}
}

\date{}

% ============================================================
\begin{document}
% ============================================================

\maketitle
\thispagestyle{empty}

\begin{center}
\small\textbf{Keywords:} LINE-1; Retrotransposons; Neurodevelopmental disorders; Promoter regions; Purifying selection; Pleiotropy; Epigenetics; Seizure; Attention-deficit/hyperactivity disorder; Genome-wide analysis
\end{center}

\newpage

% ============================================================
\begin{abstract}
LINE-1 retrotransposons constitute approximately 17\% of the
human genome and are capable of influencing gene expression
when inserted in proximity to regulatory regions.
Neurodevelopmental disorder (NDD) genes require precise
spatiotemporal regulation during brain development;
however, the relationship between LINE-1 occupancy and the promoter
architecture of NDD-associated loci has not been systematically
examined. Here, we performed a genome-wide computational
analysis comparing LINE-1 element density in promoter regions
($\pm$2 kb from the transcription start site) across five NDD
gene sets --- including genes annotated to autism spectrum
disorder (SFARI Gene database), seizure (HPO:HP:0001250),
attention-deficit/hyperactivity disorder (HPO:HP:0007018),
their phenotypic intersection, and syndromic NDD genes ---
against a curated housekeeping gene set ($n$=2,031) derived
from the HRT Atlas. All NDD gene sets exhibited significantly
lower LINE-1 promoter occupancy compared to housekeeping genes
(Mann-Whitney U test, $p < 0.05$ for all except syndromic NDD).
A consistent gradient was observed, with genes annotated to
both seizure and ADHD phenotypes showing the lowest LINE-1
occupancy (23.5\% vs 31.1\% in housekeeping genes; $p = 0.003$,
rank-biserial $r = 0.081$). Intronic LINE-1 enrichment observed
in NDD genes was attributable to gene length confounding, as
demonstrated by length-matched control analysis. These findings
suggest that NDD gene promoters are subject to purifying
selection against LINE-1 insertions, potentially reflecting
the intolerance of these loci to transcriptional dysregulation.
The pronounced depletion in genes annotated to both seizure
and ADHD phenotypes provides a genomic context for understanding
the regulatory vulnerability of pleiotropic neurodevelopmental
loci, with implications for interpreting non-coding variation in clinical genomics.
\end{abstract}

\newpage

% ============================================================
\tableofcontents
\newpage
% ============================================================

% ============================================================
\section*{Introduction}
\addcontentsline{toc}{section}{Introduction}
% ============================================================

Long Interspersed Element-1 (LINE-1, L1) retrotransposons are
the most abundant autonomous mobile genetic elements in the
human genome, comprising approximately 17\% of total genomic
sequence\cite{ardeljan2017, beck2010}.
Full-length LINE-1 elements span approximately 6 kilobases and encode two open
reading frames: ORF1p, an RNA-binding protein, and ORF2p, a
bifunctional protein with endonuclease and reverse transcriptase
activities that together execute retrotransposition via a
target-primed reverse transcription mechanism\cite{ardeljan2017}.
Although the vast majority of genomic LINE-1 copies are
retrotransposition-incompetent due to 5$'$ truncation or
accumulated mutations, an estimated 80--100 loci per haploid
genome retain retrotransposition competence, and ongoing LINE-1
activity continues to be a source of inter-individual genomic
variation\cite{beck2010}.

Beyond their role as mobile elements, LINE-1 sequences exert
regulatory influence on the transcriptome through multiple
mechanisms. The LINE-1 5$'$ untranslated region harbors
internal sense promoter activity, the strength of which is
substantially modulated by upstream flanking genomic sequences
at the insertion site\cite{lavie2004}. Additionally, the same
region contains an antisense promoter capable of initiating
transcription of neighboring sequences and influencing the
expression of adjacent protein-coding genes\cite{honda2020}.
In human neural progenitor cells, genome-wide loss of CpG
methylation leads to selective transcriptional activation of
evolutionarily young, hominoid-specific LINE-1 elements, which
subsequently act as alternative promoters for hundreds of
protein-coding genes enriched for neuronal functions\cite{jonsson2019}.
Notably, several genes regulated through this LINE-1--mediated
promoter mechanism have been independently implicated in
neurodevelopmental disorders, including autism, cognitive
impairment, and epilepsy\cite{jonsson2019}.

LINE-1 elements exhibit elevated activity during neuronal
development and are subject to relaxed silencing constraints
in neural lineages relative to other somatic tissues.
Muotri et al. demonstrated that engineered human LINE-1 elements can
retrotranspose in rat hippocampal neural progenitor cells in
vitro and produce somatic neuronal mosaicism in transgenic mice
in vivo\cite{muotri2005}.
Coufal et al. subsequently confirmed
that human LINE-1 elements retrotranspose in neural progenitors
derived from human fetal brain and embryonic stem cells, and
detected a statistically significant increase in endogenous
LINE-1 copy number in hippocampal tissue relative to non-neural
somatic tissues from the same individuals\cite{coufal2009}.
These findings indicate that LINE-1 retrotransposition occurs
preferentially in the neural lineage and may contribute to
somatic genomic diversity among post-mitotic neurons\cite{singer2010}.
The preferential activity of LINE-1 in neural progenitors is
mechanistically linked to developmentally regulated epigenetic
remodeling: LINE-1 promoter methylation decreases transiently
during the transition from neural stem cells to progenitors,
coinciding with a window of permissive retrotransposition
activity\cite{muotri2005, jonsson2019}.

The proximity of LINE-1 elements to gene regulatory regions
raises the question of whether genomic loci under strong
functional constraint exhibit systematic differences in LINE-1
occupancy relative to less constrained regions. Genes
associated with neurodevelopmental disorders (NDDs) are
characterized by unusually strict spatiotemporal regulation
during brain development, haploinsufficiency intolerance, and
high constraint scores in population genetics databases.
If LINE-1 insertion into the promoter regions of NDD-associated
genes is deleterious and subject to negative selection, these
loci would be expected to show lower LINE-1 occupancy relative
to constitutively expressed housekeeping genes, which are under
comparatively relaxed regulatory constraints.
To our knowledge, this prediction has not been systematically examined at genome
scale across multiple NDD gene sets.

Here, we performed a genome-wide computational analysis comparing LINE-1 element
density in promoter regions ($\pm$2 kb from the transcription
start site) across five independently curated NDD gene sets and
a housekeeping gene control set derived from the HRT Atlas,
using pre-computed RepeatMasker annotations for the GRCh38
assembly and intersection analysis with BEDTools.
We further assessed whether LINE-1 occupancy patterns extend to intronic
regions and whether observed differences are confounded by
gene length.

% ============================================================
\section*{Materials and Methods}
\addcontentsline{toc}{section}{Materials and Methods}
% ============================================================

\subsection*{Data Sources and Gene Set Curation}
\addcontentsline{toc}{subsection}{Data Sources and Gene Set Curation}

Five NDD-associated gene sets were curated from two independent,
publicly available databases.
Autism spectrum disorder (ASD) candidate genes were obtained
from the SFARI Gene database (March 2026 release)\cite{sfari}.
Genes were stratified by evidence score: a high-confidence set
(Tier 1) comprised genes with SFARI scores 1 or 2 ($n$=952),
and an extended set (Tier 2) additionally included score 3 genes
($n$=1,173). Genes annotated as syndromic (syndromic flag = 1)
were retained as a separate subgroup ($n$=170 promoters) for
subgroup analysis. The Tier 1 set was used as the primary NDD
gene set in all main analyses.

Seizure-associated genes (HPO:HP:0001250, $n$=2,110) and
attention-deficit/hyperactivity disorder-associated genes
(HPO:HP:0007018, $n$=441) were downloaded from the Human
Phenotype Ontology (HPO) database\cite{hpo}.
A phenotypic intersection set was defined as genes annotated to both HPO
terms ($n$=346), representing loci implicated in
epilepsy--ADHD comorbidity at the annotation level. 
The phenotypic intersection of seizure and ADHD was specifically selected for analysis because it represents one of the most highly prevalent clinical comorbidities in neurodevelopment, providing a robust model to test the regulatory constraints on pleiotropic loci\cite{fan2023}.

The housekeeping gene reference set was obtained from the
HRT Atlas v1.0\cite{hounkpe2021}, which defines housekeeping
genes as transcripts stably expressed across 52 human tissue
and cell types.
The raw set comprised 2,176 unique gene
symbols. To ensure independence between gene sets, genes
present in the SFARI Tier 2 set were excluded from the
housekeeping list prior to analysis, yielding a final
housekeeping control set of 2,031 genes.

\subsection*{Genomic Coordinate Processing}
\addcontentsline{toc}{subsection}{Genomic Coordinate Processing}

Genomic annotations were based on the GRCh38 (hg38) reference
assembly. Protein-coding gene coordinates were extracted from
GENCODE v47\cite{frankish2023} using the comprehensive
annotation GTF file. For each protein-coding gene, promoter
regions were defined as $\pm$2 kb windows centered on the
transcription start site (TSS), resolved in a strand-aware
manner: for genes on the positive strand, the TSS was defined
as the annotated start coordinate; for genes on the negative
strand, the TSS was defined as the annotated end coordinate.
Windows extending below coordinate zero were excluded. This
procedure yielded 20,092 unique protein-coding gene promoter
regions.

Intronic coordinates were derived from the same GTF file.
Transcript-level exon boundaries were extracted for all
protein-coding transcripts ($n$=65,049). For each transcript
with two or more exons, intronic intervals were computed as
the regions between consecutive exon end and start coordinates,
excluding the exon boundaries themselves. This yielded 658,138
unique intronic intervals across all transcripts.

All genomic interval files were formatted as BED files and
sorted by chromosome and start coordinate prior to intersection
analysis.

\subsection*{LINE-1 Annotation and Intersection Analysis}
\addcontentsline{toc}{subsection}{LINE-1 Annotation and Intersection Analysis}

LINE-1 element coordinates were obtained from the UCSC
RepeatMasker track for the hg38 assembly\cite{smit2015},
downloaded as a flat text file
(\texttt{rmsk.txt}). Elements were filtered to retain only
those classified as LINE/L1 family ($n$=1,031,524), and
reformatted as a sorted BED file.

Intersection analyses between gene set promoter regions and
LINE-1 elements were performed using BEDTools v2.31.1\cite{quinlan2010}.
For each promoter or intronic interval, the number of
overlapping LINE-1 elements was computed using the
\texttt{bedtools intersect -c} command with the \texttt{-sorted}
flag. Gene set-specific BED files were generated by filtering
the genome-wide promoter and intron coordinate files against
each curated gene list.

\subsection*{Statistical Analysis}
\addcontentsline{toc}{subsection}{Statistical Analysis}

All statistical analyses were performed in Python 3.13 using
the \texttt{scipy.stats} and \texttt{pandas} libraries.
LINE-1 element counts per promoter were compared between
each NDD gene set and the housekeeping control using the
two-sided Mann-Whitney U test, a non-parametric test
appropriate for count data with non-normal distributions.
Effect size was quantified as the rank-biserial correlation
coefficient $r$, calculated as:

\begin{equation}
r = 1 - \frac{2U}{n_1 n_2}
\end{equation}

\noindent where $U$ is the Mann-Whitney U statistic and $n_1$,
$n_2$ are the sample sizes of the two groups. A positive $r$
indicates that the NDD gene set has lower LINE-1 counts than
the housekeeping set. Approximate 95\% confidence intervals for $r$ were computed
using the normal approximation to the standard error:

\begin{equation}
\mathrm{SE}(r) \approx 1.96 \sqrt{\frac{n_1 + n_2 + 1}{3\,n_1\,n_2}}
\end{equation}

For the intronic confounding assessment, gene-level summaries
were computed by aggregating total intronic length (in kb) and
total LINE-1 count per gene. Intronic LINE-1 density was
expressed as LINE-1 elements per kilobase. A length-matched
control analysis was performed by identifying, for each NDD Tier
1 gene, the housekeeping gene with the most similar total
intronic length within a $\pm$50\% tolerance window. Each housekeeping gene was matched to at most one NDD gene; matched pairs were excluded from further matching to ensure independent observations. Mann-Whitney U test and rank-biserial
$r$ were then computed on the matched pairs. Spearman rank
correlation between total intronic length and LINE-1 density
was used to quantify the confounding effect.

A $p$-value threshold of 0.05 was applied for all comparisons.
No correction for multiple testing was applied across gene
sets, as each comparison was pre-specified and hypothesis-driven.

% ============================================================
\section*{Results}
\addcontentsline{toc}{section}{Results}
% ============================================================

\subsection*{LINE-1 Occupancy Is Reduced in NDD Gene Promoters}
\addcontentsline{toc}{subsection}{LINE-1 Occupancy Is Reduced in NDD Gene Promoters}

Genome-wide intersection of 1,031,524 LINE-1 elements with
promoter regions ($\pm$2 kb from TSS) of 20,092 protein-coding
genes revealed substantial variability in LINE-1 occupancy
across functional gene categories. Among housekeeping genes
($n$=1,982 promoters), 31.1\% harbored at least one LINE-1
element, establishing the baseline occupancy rate for
constitutively expressed loci.

The primary NDD gene set (SFARI Tier 1, $n$=945 promoters)
exhibited significantly lower LINE-1 occupancy compared to
housekeeping genes (27.2\%; Mann-Whitney U test: $p = 0.014$,
rank-biserial $r = 0.045$; \Tab{tab:main}).
This depletion was
consistent across gene sets curated from an independent
annotation resource: genes annotated to the seizure phenotype
in the Human Phenotype Ontology (HPO:HP:0001250, $n$=2,072
promoters) showed 28.6\% occupancy ($p = 0.037$, $r = 0.031$),
and genes annotated to attention-deficit/hyperactivity
disorder (HPO:HP:0007018, $n$=435 promoters) showed 24.8\%
occupancy ($p = 0.005$, $r = 0.069$; \Tab{tab:main}).
These results indicate that LINE-1 depletion from promoter
regions is a consistent feature of NDD-associated loci,
independent of the annotation database used to define the
gene sets.

\subsection*{A Gradient of LINE-1 Depletion Across NDD Subgroups}
\addcontentsline{toc}{subsection}{A Gradient of LINE-1 Depletion Across NDD Subgroups}

Subgroup analyses revealed a monotonic gradient of LINE-1
depletion that intensified with increasing phenotypic
specificity and comorbidity burden (\Fig{fig:main}).
Syndromic NDD genes ($n$=170 promoters), defined as SFARI
genes with confirmed syndromic presentation, showed 24.7\%
LINE-1 occupancy. Although the difference from housekeeping
genes did not reach statistical significance at the 0.05
threshold ($p = 0.060$, $r = 0.071$), likely attributable to
the reduced sample size, the direction of effect was
consistent with the broader pattern.

The most pronounced depletion was observed in the phenotypic
intersection set --- genes annotated to both HPO:HP:0001250
(seizure) and HPO:HP:0007018 (ADHD) in the HPO database
($n$=340 promoters). This set exhibited only 23.5\% LINE-1
occupancy, representing the lowest value across all gene sets
examined and a statistically significant difference from the
housekeeping baseline ($p = 0.003$, $r = 0.081$). It is important
to note that this intersection is defined purely at the level
of HPO annotation; it does not imply that all genes in this
set are causally implicated in both phenotypes, as a single
gene may be annotated to multiple phenotypes through distinct
mechanisms or alleles.

Across all six groups, LINE-1 occupancy followed a consistent
gradient: Housekeeping (31.1\%) $>$ HPO Seizure (28.6\%) $>$
NDD Tier 1 (27.2\%) $>$ HPO ADHD (24.8\%) $\approx$
Syndromic NDD (24.7\%) $>$ HPO Seizure $\cap$ ADHD (23.5\%).
All effect sizes were positive, indicating uniform directionality
of depletion (\Fig{fig:main}B).

\subsection*{Intronic LINE-1 Enrichment Is Attributable to
Gene Length Confounding}
\addcontentsline{toc}{subsection}{Intronic LINE-1 Enrichment Is Attributable to Gene Length Confounding}

To assess whether the promoter depletion pattern extended to
other genic regions, LINE-1 intersection was additionally
performed on intronic coordinates. At the intron level, NDD
Tier 1 genes exhibited higher LINE-1 occupancy than housekeeping
genes (29.8\% vs 22.7\%; $p = 4.3 \times 10^{-121}$,
$r = 0.056$). However, this apparent enrichment was driven by
a pronounced difference in gene architecture: the median total
intronic length of NDD Tier 1 genes was 256.4 kb, compared to
37.2 kb for housekeeping genes --- a 6.9-fold difference.
A Spearman rank correlation between total intronic length and
intronic LINE-1 density confirmed a significant positive
association in NDD genes ($\rho = 0.30$, $p = 7.9 \times 10^{-21}$),
indicating that longer genes accumulate more LINE-1 elements
by virtue of their size.

To directly test whether the intronic enrichment signal was
confounded by gene length, a length-matched control analysis
was performed. For each NDD Tier 1 gene ($n$=933), a
housekeeping gene with total intronic length within $\pm$50\%
was identified without replacement, yielding 678 matched pairs.
After matching, the difference in intronic LINE-1 density
between NDD and housekeeping genes was no longer statistically
significant ($p = 0.208$, $r = 0.037$), demonstrating that the
intronic signal is fully attributable to gene length
confounding rather than a genuine depletion or enrichment of
LINE-1 in NDD intronic regions.

In contrast, the promoter depletion signal was not subject to
this confounding, as all promoter windows were defined as
fixed-length 4 kb intervals centered on the TSS.
Taken together, these results indicate that LINE-1 depletion in NDD
genes is specific to promoter regulatory regions and does not
extend to intronic sequences once gene length is controlled.

% ============================================================
\section*{Discussion}
\addcontentsline{toc}{section}{Discussion}
% ============================================================

That gene promoters are depleted of transposable elements is
not a new observation. Simonti et al. showed that across
hundreds of human cell lines and primary tissues, only 5.1\%
of annotated gene promoters overlap a TE --- roughly one-tenth
of what would be expected under a random genomic
distribution\cite{simonti2017}. A complementary analysis by
You et al. found that CpG islands containing transcription
start sites are almost entirely free of LINE, SINE, and LTR
elements, and proposed that this pattern persists through
counterselection against TE insertions into these regions in
the germline\cite{you2025}. The present analysis extends
these observations in a specific direction: not all gene
promoters are equally depleted. Neurodevelopmental disorder
genes, defined across five independent annotation sources,
show consistently lower LINE-1 occupancy than housekeeping
genes, and the magnitude of that difference scales with the
regulatory demands of the associated phenotype.

The gradient from housekeeping genes through seizure, NDD,
ADHD, and finally seizure-ADHD comorbidity genes is
consistent with the idea that functional constraint on these
loci is not uniform. NDD-associated genes, particularly
those encoding transcriptional regulators, are characterized
by extreme dosage sensitivity: a 50\% reduction in gene
product is sufficient to produce a developmental
phenotype\cite{zug2022}. Multiple autism-associated
transcription factors have been shown to converge on a
narrow window 1--2 kb upstream of the TSS, with this
promoter-proximal overlap being highly predictive of
brain-expressed genes\cite{darbandi2024}. The same region
has been shown to carry open chromatin and active histone
marks specifically in neural progenitor cells for ASD-risk
genes\cite{koesterich2023}. A de novo LINE-1 insertion landing in this zone would be recognized by the host epigenetic machinery, recruiting TRIM28/KAP1 and leading to H3K9me3 deposition, repressive marks that can spread into flanking sequence\cite{jonsson2019} (\Fig{fig:model}). Beyond local chromatin effects, Lu et al. showed that LINE-1-enriched genomic regions are
physically sequestered in nucleolus- and lamina-associated
inactive nuclear compartments, while SINE-enriched
housekeeping genes reside in the transcriptionally active
nuclear interior\cite{lu2020}. A de novo LINE-1 insertion
at the promoter of an NDD gene could therefore disrupt not
only local transcription factor access but also the broader
nuclear positioning that maintains its expression.

The depletion signal is region-specific. Intronic LINE-1
differences between NDD and housekeeping genes disappeared
after length correction, indicating that the extended gene
bodies of NDD loci simply accumulate more elements over time,
as expected. Interspersed repeats in general tend to
accumulate at distances from the TSS rather than in
immediately proximal promoter regions\cite{liang2015}, which
makes the promoter window the appropriate place to detect any
functional depletion. The fact that the signal survives only
in this fixed regulatory interval, and not in introns once
gene length is matched, is consistent with selection acting
specifically on the transcription-initiating region rather
than on the gene body as a whole.

The observation that genes annotated to both seizure and ADHD
phenotypes show the deepest depletion could reflect the
cumulative regulatory burden on pleiotropic loci. Genes
implicated in two distinct neurodevelopmental phenotypes may
face constraints from two sets of cell-type-specific
regulatory programs, making their promoter regions doubly
intolerant to disruptive insertions. It should be noted,
however, that this intersection set is defined by HPO
annotation overlap rather than confirmed comorbidity
genetics, so this interpretation is speculative and would
require independent validation with curated comorbidity
gene lists.

Taken together, the pattern is consistent with the hypothesis
that functional constraint at NDD gene promoters has, over
primate evolutionary time, selected against LINE-1 insertions in
these regions. The small effect sizes observed ($r = 0.031$--0.081) are expected: this would be a weak but
consistent force acting on a large number of loci. Genome-wide
patterns of TE depletion at regulatory regions reflect the
aggregate outcome of many independent insertion and selection
events across deep evolutionary time, and individually small
per-gene signals are characteristic of such
genome-wide constraint analyses\cite{simonti2017,
chuong2016}. Whether the depletion genuinely reflects
germline selection or is partly explained by local sequence
composition differences (e.g., GC content disparities between NDD and housekeeping promoters) cannot be definitively resolved from the present genomic snapshot alone.

\subsection*{Limitations}
\addcontentsline{toc}{subsection}{Limitations}

Several limitations of this study warrant consideration.
First, the analysis is based on binary LINE-1 occupancy
(presence or absence per promoter window) and does not account
for LINE-1 element age, subfamily composition, or degree of
truncation, all of which may influence regulatory impact.
Second, all gene sets were defined at the level of database
annotation, and individual genes within these sets may vary
substantially in their biological roles and phenotypic
relationships. In particular, the HPO Seizure $\cap$ ADHD
intersection does not represent a curated comorbidity gene
set, but rather the overlap of two independent annotation
lists. Third, the promoter window size of $\pm$2 kb was
selected as a standard regulatory interval; the depletion
pattern may differ with alternative window definitions.
Fourth, effect sizes across all comparisons were small
(rank-biserial $r = 0.031$--0.081), indicating that while
the signal is consistent and statistically robust, the
biological magnitude of LINE-1 depletion at the individual
gene level is modest. Finally, this is a purely computational
study based on a single reference genome snapshot; the
interpretation that depletion reflects historical negative
selection against LINE-1 insertions at NDD promoters would
require population-level insertion polymorphism data and
evolutionary analyses beyond the scope of this work.

% ============================================================
\section*{Conclusions}
\addcontentsline{toc}{section}{Conclusions}
% ============================================================

This analysis shows that the promoter regions of
neurodevelopmental disorder-associated genes are consistently
depleted of LINE-1 retrotransposons relative to housekeeping
gene promoters. The depletion is statistically significant
across multiple independently curated NDD gene sets and
follows a gradient that intensifies with phenotypic specificity,
with the strongest signal observed in genes annotated to both
seizure and ADHD phenotypes in the Human Phenotype Ontology.
Intronic LINE-1 differences between NDD and housekeeping genes
were largely explained by gene length and did not persist after
length-matched control analysis.

The region-specific pattern --- depletion at fixed-length
promoter windows but no signal in introns after length
correction --- is consistent with functional constraint acting
on LINE-1 insertions at NDD gene regulatory regions rather than
a general genomic property of these genes. What this
analysis establishes is the pattern itself: across five
independently curated gene sets, NDD promoters carry fewer
LINE-1 elements than housekeeping promoters, and the magnitude
of depletion tracks with the regulatory demands of the
phenotype. Future work combining this approach with population-level insertion polymorphism data (e.g., leveraging structural variant cohorts like gnomAD SV/MEI) could explicitly test whether this depletion represents an ongoing evolutionary signal or a neutral consequence of local sequence composition.

% ============================================================
\section*{Data and Code Availability}
\addcontentsline{toc}{section}{Data and Code Availability}
% ============================================================

All underlying datasets used in this study are publicly available. Gene set annotations were obtained from the SFARI Gene database (March 2026 release; \url{https://gene.sfari.org}) and the Human Phenotype Ontology database (\url{https://hpo.jax.org}). Genomic coordinates were sourced from GENCODE v47 (\url{https://www.gencodegenes.org}) and the UCSC RepeatMasker track for hg38 (\url{https://hgdownload.soe.ucsc.edu}). The housekeeping gene reference set was obtained from the HRT Atlas v1.0 (\url{https://housekeeping.unicamp.br}). All Python and Bash scripts used for data processing, BEDTools intersection, and statistical analysis are freely available in a dedicated GitHub repository at \url{https://github.com/cansevilmis/LINE1_NDD_Promoters}.

% ============================================================
\section*{Ethics Statement}
\addcontentsline{toc}{section}{Ethics Statement}
% ============================================================

Not applicable. This study was conducted exclusively using publicly available, anonymized genomic and phenotypic databases. No human or animal subjects were involved in this research.

% ============================================================
\section*{Author Contributions}
\addcontentsline{toc}{section}{Author Contributions}
% ============================================================

\textbf{Can Sevilmiş}: Conceptualization, Methodology, Software, Formal Analysis, Data Curation, Writing - Original Draft, Writing - Review \& Editing, Visualization.

% ============================================================
\section*{Funding}
\addcontentsline{toc}{section}{Funding}
% ============================================================

This research received no specific grant from any funding agency in the public, commercial, or not-for-profit sectors.

% ============================================================
\section*{Acknowledgements}
\addcontentsline{toc}{section}{Acknowledgements}
% ============================================================

The author thanks the developers of SFARI Gene, Human Phenotype Ontology, GENCODE, UCSC Genome Browser, HRT Atlas, and BEDTools for providing open-access resources and software essential to this work.

% ============================================================
\section*{Declaration of Generative AI and AI-Assisted Technologies}
% ============================================================

During the preparation of this manuscript, the author used Claude (Anthropic) to assist with writing and troubleshooting bioinformatic scripts, as well as for grammatical and stylistic refinement of the text. The author thoroughly reviewed and edited the content following the use of these tools and takes full responsibility for the final content of the publication.

% ============================================================
\section*{Competing Interests}
% ============================================================

The author declares no competing interests.

% ============================================================
\FloatBarrier

\begin{figure}[htbp]
\centering
\includegraphics[width=0.99\textwidth]{final_panel_v2.png}
\caption{\textbf{LINE-1 retrotransposon occupancy in promoter
regions across neurodevelopmental disorder gene sets.}
(A) Proportion of promoters ($\pm$2 kb from TSS) containing
at least one LINE-1 element across five NDD gene sets and a
housekeeping gene control.
The dashed line indicates the
housekeeping baseline (31.1\%). Significance markers above
bars denote Mann-Whitney U test results versus the housekeeping
gene set (* $p < 0.05$, ** $p < 0.01$, n.s. not significant).
(B) Effect sizes (rank-biserial $r$) with approximate 95\%
confidence intervals for each NDD gene set compared to the
housekeeping gene set.
All point estimates are positive and
their confidence intervals do not cross zero for statistically
significant comparisons, indicating consistent LINE-1 depletion
relative to housekeeping genes across all NDD groups.
Gene set sample sizes are indicated in parentheses on the
x-axis (A) and y-axis labels (B).
NDD Tier 1: SFARI Gene
database, scores 1--2. Syndromic NDD: SFARI genes with
syndromic annotation. HPO Seizure: HPO:HP:0001250.
HPO ADHD: HPO:HP:0007018.
HPO Seizure $\cap$ ADHD: genes annotated to both HPO terms.}
\label{fig:main}
\end{figure}

% ============================================================
\FloatBarrier

\begin{figure}[htbp]
\centering
\includegraphics[width=0.95\textwidth]{line1_silencing_diagram.png}
\caption{\textbf{Conceptual model of LINE-1 mediated epigenetic silencing at NDD gene promoters.} 
(Top) Housekeeping gene promoters, which are depleted of transposable elements, maintain an open chromatin architecture (euchromatin) that readily permits the binding of transcription factors and RNA Polymerase II at the transcription start site (TSS). 
(Bottom) De novo insertion of a LINE-1 retrotransposon into the proximal promoter region of a highly constrained neurodevelopmental disorder (NDD) gene is recognized by host silencing machinery, recruiting TRIM28/KAP1. This initiates the local deposition of repressive H3K9me3 histone marks and subsequent chromatin condensation (heterochromatin), physically occluding the TSS and leading to transcriptional repression.}
\label{fig:model}
\end{figure}

% ============================================================
\FloatBarrier

\begin{table}[htbp]
\centering
\caption{\textbf{Summary of LINE-1 promoter occupancy analysis
across gene sets.} All comparisons were performed against the
housekeeping gene set ($n$=1,982) using the two-sided
Mann-Whitney U test.
Effect size is reported as the
rank-biserial correlation coefficient ($r$); positive values
indicate lower LINE-1 counts in the NDD gene set relative
to housekeeping genes.
Approximate 95\% confidence intervals
for $r$ are provided in parentheses. n.s.
= not significant
($p \geq 0.05$).}
\label{tab:main}
\small
\begin{tabularx}{\textwidth}{>{\raggedright\arraybackslash}Xcccc}
\toprule
\textbf{Gene Set} &
\textbf{\textit{n}} &
\textbf{LINE-1 (\%)} &
\textbf{\textit{p}-value} &
\textbf{\textit{r} (95\% CI)} \\
\midrule
Housekeeping
  & 1,982 & 31.1 & --- & --- \\
HPO Seizure
  & 2,072 & 28.6 & 0.037 & 0.031 (0.00, 0.07) \\
NDD Tier 1 (SFARI score 1--2)
  &   945 & 27.2 & 0.014 & 0.045 (0.01, 0.08) \\
HPO ADHD
  &   435 & 24.8 & 0.005 & 0.069 (0.01, 0.12) \\
Syndromic NDD
  &   170 & 24.7 & 0.060 (n.s.) & 0.071 ($-$0.02, 0.16) \\
HPO Seizure $\cap$ ADHD
  &   340 & 23.5 & 0.003 & 
0.081 (0.02, 0.14) \\
\bottomrule
\end{tabularx}
\normalsize
\end{table}

% ============================================================
\footnotesize
% BibTeX kullanımı ve boşlukların sıfırlanması
\bibliographystyle{unsrtnat}
\setlength{\bibsep}{0pt}
\bibliography{references}

\end{document}